\documentclass[14pt,a4paper]{extarticle}
\usepackage{fontspec}
\usepackage{polyglossia}
\setmainlanguage{russian}
\setmainfont{Times New Roman}
\setmonofont{Courier New}
\usepackage{setspace}
\onehalfspacing
\usepackage{geometry}
\geometry{left=30mm, right=15mm, top=20mm, bottom=20mm}
\usepackage{indentfirst}
\setlength{\parindent}{1.25cm}
\setlength{\parskip}{0pt}
\usepackage{xurl}
\urlstyle{same}
\usepackage{hyperref}
\usepackage{booktabs}
\usepackage{ragged2e}
\usepackage{graphicx}
\usepackage{float}
\usepackage{pdflscape}
\usepackage{listings}
\usepackage{xcolor}
\usepackage{array}
\usepackage{amsmath} % <-- ИСПРАВЛЕНО: Подключен ДО wasysym
\usepackage{wasysym}

\lstset{
	language=C++,
	basicstyle=\small,
	keywordstyle=\color{blue},
	stringstyle=\color{red},
	commentstyle=\color{green!50!black},
	showstringspaces=false,
	breaklines=true,
	frame=single,
	numbers=left,
	numberstyle=\tiny\color{gray},
	stepnumber=1,
	keepspaces=true,
	extendedchars=true,
	inputencoding=utf8
}

\begin{document}
	
	\sloppy
	\thispagestyle{empty}
	\centering{
		Министерство образования и науки Российской Федерации\\
		Федеральное государственное автономное образовательное учреждение высшего образования\\
		«Санкт-Петербургский политехнический университет Петра Великого»\\[10pt]
		Институт компьютерных наук и кибербезопасности\\
		Высшая школа технологий искусственного интеллекта\\
		02.03.01 Математика и компьютерные науки\\
		\vspace{\fill}
		{\large
			Отчёт по дисциплине\\
			{\Large \bfseries Вычислительная математика}\\
			Лабораторная работа №2\\
			{\Large <<Аппроксимация>>}\\
			Вариант 18\\
			\vspace{\fill}
		}
		
		{
			\flushright
			{\bfseries Выполнил:}\\
			студент группы 5130201/40003\\
			Джабраилов А.В.\\
			\hspace{9cm}Подпись: \hrulefill \\
			{\bfseries Принял:}\\
			К. физ.-мат. наук, доц., доц.\\
			Пак В.Г.\\
			\hspace{9cm}Подпись: \hrulefill \\
			
			\hspace{9cm}<<\rule{1.5cm}{0.4pt}>> \hrulefill\ 2026~г.
		}
		
		\centering{
			Санкт-Петербург, 2026
		}
	}
	
	\justifying
	
	\newpage
	\sloppy
	
	\section*{1 Задание}
	Для данной табличной функции построить многочлены с наилучшим среднеквадратичным отклонением первой и второй степени методом наименьших квадратов.
	
	\section*{2 Цель работы}
	Научиться аппроксимировать табличные функции полиномами методом наименьших квадратов.
	
	\section*{3 Ход работы}
	
	\begin{table}[h]
		\centering
		\caption{Таблица значений (часть 1).}
		\begin{tabular}{|c|c|c|c|c|c|c|c|c|c|c|}
			\hline
			$x_i$ & 0,1 & 0,3 & 0,5 & 0,7 & 0,9 & 1,1 & 1,3 & 1,5 & 1,7 & 1,9 \\ \hline
			$y_i$ & -0,86 & -0,77 & -0,56 & -0,46 & -0,28 & -0,24 & -0,36 & -0,43 & -0,56 & -0,59 \\ \hline
		\end{tabular}
	\end{table}
	
	\begin{table}[h]
		\centering
		\caption{Таблица значений (часть 2).}
		\begin{tabular}{|c|c|c|c|c|c|c|c|c|c|c|}
			\hline
			$x_i$ & 2,1 & 2,3 & 2,5 & 2,7 & 2,9 & 3,1 & 3,3 & 3,5 & 3,7 & 3,9 \\ \hline
			$y_i$ & -0,70 & -1,01 & -1,03 & -1,47 & -1,68 & -1,93 & -2,28 & -2,53 & -2,93 & -3,07 \\ \hline
		\end{tabular}
	\end{table}
	
	Алгоритм вычислений следующий:
	\begin{enumerate}
		\item Построение матрицы Вандермонда, т.к. т.о. можно пробразовать задачу подбора коэффициентов многочлена в задачу решения СЛУ: умножением матрицы $A$ на вектор неизвестных коэффициентов $a$ получаем вектор значений многочлена в точках $x$;
		\item Вычисление коэффициентов $a$ методом Гаусса;
		\item Вычисление среднеквадратичного отклонения.
	\end{enumerate}
	
	Для выполнения работы была написана программа на языке \textit{C++}. Программа имеет модульную структуру и состоит из четырёх функций:
	\begin{enumerate}
		\item \textit{solveLinearSystem} — решает систему линейных уравнений $Ax=b$.
		\item \textit{polyval} — вычисляет значение полинома.
		\item \textit{approximation} — реализует метод наименьших квадратов.
		\item \textit{main} — точка входа программы.
	\end{enumerate}
	
	Для вычислений используются библиотеки \textit{vector} и \textit{cmath}. 
	
	Текст исходного кода функций приведён в приложении.
	
	\section*{4 Результаты}
	\begin{enumerate}
		\item Многочлен первой степени: $P_1(x) = 0.13 - 0.66x$. Среднее квадратичное отклонение: 0,45.
		\item Многочлен второй степени: $P_2(x) = -0.87 + 0.84x - 0.38x^2$. Среднее квадратичное отклонение: 0,10.
	\end{enumerate}
	
	\section*{5 Приложение}
	
	\begin{lstlisting}[caption={Файл Lab2.cpp}, label={lst:fullcode}]
#include <iostream>
#include <vector>
#include <cmath>
#include <iomanip>

// Solve linear system Ax = b using Gaussian elimination with partial pivoting
std::vector<double> solveLinearSystem(std::vector<std::vector<double>> A, std::vector<double> b) {
	int n = b.size();
	
	// Forward elimination with partial pivoting
	for (int k = 0; k < n; k++) {
		// Find pivot
		int maxRow = k;
		for (int i = k + 1; i < n; i++) {
			if (std::abs(A[i][k]) > std::abs(A[maxRow][k])) {
				maxRow = i;
			}
		}
		
		// Swap rows
		std::swap(A[k], A[maxRow]);
		std::swap(b[k], b[maxRow]);
		
		// Eliminate column
		for (int i = k + 1; i < n; i++) {
			double factor = A[i][k] / A[k][k];
			for (int j = k; j < n; j++) {
				A[i][j] -= factor * A[k][j];
			}
			b[i] -= factor * b[k];
		}
	}
	
	// Back substitution
	std::vector<double> x(n);
	for (int i = n - 1; i >= 0; i--) {
		x[i] = b[i];
		for (int j = i + 1; j < n; j++) {
			x[i] -= A[i][j] * x[j];
		}
		x[i] /= A[i][i];
	}
	
	return x;
}

// Evaluate polynomial at point x
double polyval(const std::vector<double>& coeffs, double x) {
	double result = 0.0;
	for (int i = coeffs.size() - 1; i >= 0; i--) {
		result = result * x + coeffs[i];
	}
	return result;
}

// Polynomial approximation using least squares
std::pair<std::vector<double>, double> approximation(
const std::vector<double>& x, 
const std::vector<double>& y, 
int degree) {
	
	int n = x.size();
	int m = degree + 1;
	
	// Build Vandermonde matrix (increasing powers)
	std::vector<std::vector<double>> A(n, std::vector<double>(m));
	for (int i = 0; i < n; i++) {
		for (int j = 0; j < m; j++) {
			A[i][j] = std::pow(x[i], j);
		}
	}
	
	// Solve normal equations: (A^T * A) * coeffs = A^T * y
	std::vector<std::vector<double>> AtA(m, std::vector<double>(m, 0.0));
	std::vector<double> Aty(m, 0.0);
	
	for (int i = 0; i < m; i++) {
		for (int j = 0; j < m; j++) {
			for (int k = 0; k < n; k++) {
				AtA[i][j] += A[k][i] * A[k][j];
			}
		}
		for (int k = 0; k < n; k++) {
			Aty[i] += A[k][i] * y[k];
		}
	}
	
	// Solve for coefficients
	std::vector<double> coeffs = solveLinearSystem(AtA, Aty);
	
	// Calculate fitted values and MSE
	double mse = 0.0;
	for (int i = 0; i < n; i++) {
		double y_fit = polyval(coeffs, x[i]);
		mse += std::pow(y[i] - y_fit, 2);
	}
	mse /= n;
	
	double rmse = std::sqrt(mse);
	
	return {coeffs, rmse};
}

int main() {
	std::vector<double> x = {
		0.1, 0.3, 0.5, 0.7, 0.9, 1.1, 1.3, 1.5, 1.7, 1.9,
		2.1, 2.3, 2.5, 2.7, 2.9, 3.1, 3.3, 3.5, 3.7, 3.9
	};
	
	std::vector<double> y = {
		-0.86, -0.77, -0.56, -0.46, -0.28, -0.24, -0.36, -0.43, -0.56, -0.59,
		-0.70, -1.01, -1.03, -1.47, -1.68, -1.93, -2.28, -2.53, -2.93, -3.07
	};
	
	// Perform approximations
	auto [coeffs_1, rmse_1] = approximation(x, y, 1);
	auto [coeffs_2, rmse_2] = approximation(x, y, 2);
	
	// Output results
	std::cout << std::fixed << std::setprecision(2);
	
	std::cout << "\nFirst-degree polynomial (linear approximation):" << std::endl;
	std::cout << "P_1(x) = " << coeffs_1[0] << " + " << coeffs_1[1] << "*x" << std::endl;
	std::cout << "Root mean square error: " << rmse_1 << std::endl;
	
	std::cout << "\nSecond-degree polynomial (quadratic approximation):" << std::endl;
	std::cout << "P_2(x) = " << coeffs_2[0] << " + " << coeffs_2[1] << "*x + "
	<< coeffs_2[2] << "*x^2" << std::endl;
	std::cout << "Root mean square error: " << rmse_2 << std::endl;
	
	return 0;
}
	\end{lstlisting}
	
\end{document}